% GENERADO por generar.py desde tesis.yaml. No editar a mano.
\begin{landscape}
\anexo{Anexo 1. Matriz de consistencia}
\noindent{\small\setstretch{1}\textbf{Título:} SISTEMA WEB CON CHATBOT DE WHATSAPP E INTELIGENCIA ARTIFICIAL PARA LA GESTIÓN DE ATENCIÓN AL CLIENTE EN UNA EMPRESA DE SERVICIOS\newline \textbf{Autor:} Apellido Apellido, Nombre Nombre\par}
{\small\setlength{\tabcolsep}{4pt}
\begin{longtable}{L{3.9cm}L{3.9cm}L{3.9cm}L{4.4cm}L{5.4cm}}
\caption{Matriz de consistencia}\label{tab:matriz-consistencia}\\
\toprule
\textbf{Problemas} & \textbf{Objetivos} & \textbf{Hipótesis} & \textbf{Variables e indicadores} & \textbf{Metodología} \\
\midrule
\endfirsthead
\multicolumn{5}{@{}l}{\itshape (continuación)}\\
\toprule
\textbf{Problemas} & \textbf{Objetivos} & \textbf{Hipótesis} & \textbf{Variables e indicadores} & \textbf{Metodología} \\
\midrule
\endhead
\bottomrule
\endlastfoot
\textbf{Problema general:}\newline ¿En qué medida un sistema web con chatbot de WhatsApp e inteligencia artificial mejora la gestión de atención al cliente en una empresa de servicios? & \textbf{Objetivo general:}\newline Determinar en qué medida un sistema web con chatbot de WhatsApp e inteligencia artificial mejora la gestión de atención al cliente en una empresa de servicios. & \textbf{Hipótesis general:}\newline Un sistema web con chatbot de WhatsApp e inteligencia artificial mejora significativamente la gestión de atención al cliente en una empresa de servicios. & \textbf{VI:} Sistema web con chatbot de WhatsApp e inteligencia artificial\newline \textbf{VD:} Gestión de atención al cliente\newline \textbf{Dimensiones:} Oportunidad de la atención; Eficacia de la atención & \textbf{Enfoque:} Cuantitativo\newline \textbf{Tipo:} Aplicada\newline \textbf{Nivel:} Explicativo\newline \textbf{Diseño:} Pre-experimental con pre-prueba y post-prueba\newline \textbf{Esquema:} GE: O1 X O2\newline \textbf{Método:} Deductivo\newline \textbf{Población:} EJEMPLO: consultas recibidas por WhatsApp en un mes típico (N = 1200).\newline \textbf{Muestra:} EJEMPLO: 30 días de registro en pre-prueba y 30 días en post-prueba.\newline \textbf{Muestreo:} No probabilístico por conveniencia\newline \textbf{Técnicas:} Observación\newline \textbf{Instrumentos:} Ficha de registro \\
\midrule
\textbf{PE1:} ¿En qué medida un sistema web con chatbot de WhatsApp e inteligencia artificial reduce el tiempo promedio de primera respuesta al cliente? & \textbf{OE1:} Determinar en qué medida un sistema web con chatbot de WhatsApp e inteligencia artificial reduce el tiempo promedio de primera respuesta al cliente. & \textbf{HE1:} Un sistema web con chatbot de WhatsApp e inteligencia artificial reduce significativamente el tiempo promedio de primera respuesta al cliente. & \textbf{Indicador:} Tiempo promedio de primera respuesta (TPR) &  \\
\midrule
\textbf{PE2:} ¿En qué medida un sistema web con chatbot de WhatsApp e inteligencia artificial incrementa la tasa de consultas resueltas sin intervención humana? & \textbf{OE2:} Determinar en qué medida un sistema web con chatbot de WhatsApp e inteligencia artificial incrementa la tasa de consultas resueltas sin intervención humana. & \textbf{HE2:} Un sistema web con chatbot de WhatsApp e inteligencia artificial incrementa significativamente la tasa de consultas resueltas sin intervención humana. & \textbf{Indicador:} Tasa de consultas resueltas sin intervención humana (TRA) &  \\
\end{longtable}}
\vspace{-10pt}
\nota{Elaboración propia.}
\end{landscape}
